\scanpage{21}
\mspage{10}
\uline{Remarque} Quand $\pi\neq(1)$, \inserted{i.e. $K$ pas fini sur $\mathbb Q$,} le théorème 1 peut se renforcer, \inserted{sauf erreur,} en établissant que pour \unsure{tt|tout} ss-groupe \inserted{$\pi'\subset\pi$} ouvert \uline{dans $\pi$}, $\operatorname{Centr}_{\unsure{\mathcal E|E}}(\pi')=\{1\}$. Si $\unsure{Z|\mathcal Z}$ est ce centralisateur, on a $\unsure{Z|\mathcal Z}\cap\pi=(1)$ d'après le th.1, \deleted{et}\inserted{prouvons que} l'image de $\unsure{Z|\mathcal Z}$ dans $\Gamma_{\bar K/K}\subset\Gamma_{\overline{\mathbb Q}/\mathbb Q}$ est \deleted{contenue dans le \unsure{noy|noyau}}\inserted{finie} (ce qui prouvera \inserted{déjà} alors, que $\unsure{Z|\mathcal Z}$ est d'ordre $1$ ou $2$, et dans le deuxième cas que son image dans $\Gamma_{\overline{\mathbb Q}/\mathbb Q}$ est \deleted{un \unsure{él.}} engendrée par un $\tau$ de conjugaison complexe) --- Quitte à remplacer $\unsure{\mathcal E|E}$ par un sous-groupe ouvert assez petit (ce qui revient à passer à une extension finie de $K$) OPS $\pi'=\pi$, alors l'image $\unsure{Z'|\mathcal Z'}$ de $\unsure{Z|\mathcal Z}$ dans $\Gamma$ est contenue dans le noyau de l'hom. $\varphi:\Gamma\longrightarrow\operatorname{Aut\,ext}(\pi)$. Or (sauf erreur) je sais prouver que cet hom. est injectif (on est ramené aussitôt au cas où $K$ est de degré de transc. $1$, et on est ramené au cas \inserted{du $\pi_1$} d'une courbe algébrique\ldots)
\margin{20}
